\input{preambel} 	% Alle Packages und Formatierungen werden in preamble.tex eingetragen um den Text hier übersichtlicher zu gestalten.

\begin{document}

\input{title}			% Deckblatt (leicht modifizierte KM-Version)
 
\begingroup
\pagestyle{empty} % ohne Seitenzahl
\tableofcontents % Inhaltsverzeichnis
\clearpage % Seite bis zum Ende leeren
\endgroup

\newpage

\setcounter{page}{1} % Seitenzähler wird auf 0 gesetzt



 
\section{Einleitung} % Überschrift für erstes Kapitel, optional auch: [Eintrag im IHV]{Einleitung}
\label{sec:Einleitung} % Textmarke, um darauf zu verweisen.
\lipsum[2]			% Dummytext (einfach löschen und los schreiben)

\section{Thematik in der Literatur}
\lipsum[30]				% Dummytext (einfach löschen und los schreiben)



\begin{figure}[!ht]		% Bild mit Unterschrift einbinden
\centering
\includegraphics[scale=1]{pics/bild.png} 
\caption[Demobild]{Einzelne Bilder wie dieses sollten zentriert auf der Seite platziert werden. Auf eine sinnvolle Größe ist ebenso zu achten, wie auf eine konkrete Beschreibung der Abbildung.}
\label{fig:bild1}\vspace{0.8em}
\end{figure}

\lipsum[31]

\newpage
\section{Überschrift}

\renewcommand{\baselinestretch}{1.0}\normalsize 
\setlength{\columnsep}{30px} % Abstand Spalten vergrößern
\begin{multicols}{2}[\noindent \textbf{Hier könnte beispielsweise ein Verweis auf den Textauszug aus der Literatur stehen, etwa so: Der folgende Text stammt aus \citep{voss}{, S.201f}}] 

\textbf{Dies kann als Zusammenfassung für den darunter liegenden Text gesehen werden. Hin und wieder ist dieser auch fett gedruckt.}
\setcounter{linenumber}{1}
\begin{linenumbers} % Schaltet Zeilennummerierung ein
%\modulolinenumbers[5] % nur jede 5. Zeile
\lipsum[99]
\end{linenumbers}
\end{multicols} % Ende 2 Spalten
\renewcommand{\baselinestretch}{1.5}\normalsize 

\subsection{Überschrift}
\lipsum[99]

\begin{table}[!ht]
\centering
\caption[Zahlenvergleich der Werte 1 und 2 ]{Übersicht. Werte aus \citep{stall}{, S.211, Abb. 2}}\vspace{0.8em}
\label{tab:tab1}
\begin{tabular}{l|cc} \toprule
 &\multicolumn{2}{r}{universal}\\ \cmidrule{1-1}
\backslashbox{ABC}{DEF} & 1 & 2 \\ \midrule 
E & 3,34 & 4,34 \\
F & 5,45 & 6,45 \\ \bottomrule
\end{tabular}
\end{table}

\lipsum[99]

\newpage
\section{Ausblick und Zusammenfassung}
Dies ist ein Text\footnote{Dies sollte eine Fußnote werden, die den Rahmen einer Zeile sprengt. Dies lässt sich auch durch Wiederholung den Inhalts erreichen. Dies lässt sich auch durch Wiederholung den Inhalts erreichen \citep{ray}.  Dies lässt sich auch durch Wiederholung den Inhalts erreichen. \newline Dies lässt sich auch durch Wiederholung den Inhalts erreichen.}
, der auch in der Literatur (\citep{demo}{, S.22}) existiert. 
\lipsum[32]

Da pack ich doch gleich no ne Fußnote drunter.\footnote{Mal sehen, welcher Abstand hier zu vorherigen Fußnote entsteht.}

\lipsum[10]

\begin{table}[!ht]					%	Tabelle genau hier einsetzen
\centering							% Tabelle zentrieren
\caption[YYY]{Tabellenüberschrift.}\vspace{0.8em} % Tabellenbeschrift.
\label{tab:2}							% Label für den Tabellenverweis
\begin{tabular}{l|r|r} \toprule		% l = linksb., c = zentriert, r = rechtsb., | vert. Strich
  & \textbf{Spalte A} & \textbf{Spalte B} \\ \midrule 				% & ist Spaltentrennzeichen, Abschluss der Zeile mit "\\"
\textbf{Zeile 1} & Wert 1A & Wert 1B \\
\textbf{Zeile 2 }& Wert 2A & Wert 2B \\ \bottomrule		% top-, mid- & bottomrule in untersch. Strichstärken
\end{tabular}
\end{table}

\lipsum[41]


\newpage
\section{Selbständigkeitserklärung}
Hiermit erkläre ich, dass ich die vorliegende Seminararbeit selbständig verfasst und keine anderen als die angegebenen Hilfsmittel benutzt habe.\footnote{eine kleine Fußnote.}

Die Stellen der Seminararbeit, die anderen Quellen im Wortlaut oder dem Sinn nach entnommen wurden, sind durch Angaben der Herkunft kenntlich gemacht. Dies gilt auch für Zeichnungen, Skizzen, bildliche Darstellungen sowie für Quellen aus dem Internet.

\vspace{20mm}
\hrule
\vspace{5mm}
\begin{center}
\textbf{Mering}, \today
\end{center}


\newpage
\listoffigures		% Abbildungsverzeichnis erstellen

\newpage
\listoftables		% Tabellenverzeichnis erstellen

\newpage
\addcontentsline{toc}{section}{Literaturverzeichnis}
\small 						

\bibliographystyle{alphadin}
\bibliography{standard}





\newpage
\appendix
\addcontentsline{toc}{section}{Anhang}
\begin{appendix}
\section*{Anhang}
%\subsection{Anhangsteil 1}
%\lipsum[31]
%\subsection{Anhangsteil 2}
\end{appendix}
\end{document}